% --- Archon physics formalization source begin ---
% archon:physics
% archon:covers PhyXMiniProblems/problem_phyx_mini_0885.lean
% archon:source-report reports/phyx_mini/problem_phyx_mini_0885.source.json

\chapter{Physics problem phyx\_mini\_0885}
\label{ch:PhyXMiniProblems_problem_phyx_mini_0885}

\paragraph{Problem source.}
\#\# Physical scenario

High-frequency signals are often transmitted along a coaxial cable, such as the one shown in figure. For example, the cable TV hookup coming into your home is a coaxial cable. The signal is carried on a wire of radius \textbackslash{}( r\_1 \textbackslash{}) while the outer conductor of radius \textbackslash{}( r\_2 \textbackslash{}) is grounded. A soft, flexible insulating material fills the space between them, and an insulating plastic coating goes around the outside.

\#\# Question

Evaluate the inductance per meter of a cable having \textbackslash{}( r\_1 = 0.50 \textbackslash{}, \textbackslash{}text\{mm\} \textbackslash{}) and \textbackslash{}( r\_2 = 3.0 \textbackslash{}, \textbackslash{}text\{mm\} \textbackslash{}).

\#\# Answer choices

- A: A: \textbackslash{}( 0.24 \textbackslash{}

- B: B: \textbackslash{}( 0.48 \textbackslash{}

- C: C: \textbackslash{}( 0.36 \textbackslash{}

- D: D: \textbackslash{}( 0.30 \textbackslash{}

\#\# Figure caption (auxiliary; use the image as primary evidence)

The image illustrates a cylindrical coaxial cable structure. It features two main components:

1. **Inner Conductor**: 
   - Depicted as a solid cylinder at the center.
   - Labeled with "Inner conductor radius \textbackslash{}(r\_1\textbackslash{})".

2. **Outer Conductor**:
   - Shown as a hollow cylindrical shell surrounding the inner conductor.
   - Labeled with "Outer conductor".

The space between the inner and outer conductors is marked with a radius \textbackslash{}(r\_2\textbackslash{}). The diagram includes dashed lines indicating the coaxial alignment and concentric relationship between the inner and outer cylinders.

\#\# Dataset metadata

Electromagnetic Induction; Physical Model Grounding Reasoning, Numerical Reasoning

\paragraph{Recorded answer/context.}
C: C: \textbackslash{}( 0.36 \textbackslash{}

\paragraph{Figure/image path.}
/root/proposal\_for\_physic/science-mango/phyx\_mini\_run/phyx\_data/test\_image/885.png

\paragraph{Formalization target.}
create a compiling Lean file with sorry bodies at `PhyXMiniProblems/problem\_phyx\_mini\_0885.lean`. The Lean declarations must preserve the physical quantities, dimensions or dimensional roles, figure labels, governing-law hypotheses, and final relation expressed by this problem.
Use Mathlib/Physlib names found through LeanExplore where available. If a physics API is missing, introduce faithful local abstractions rather than scalar placeholder aliases.

\begin{theorem}[Physics formalization target]
\label{thm:physics:phyx_mini_0885:target}
\lean{PhyXMiniProblems.ProblemPhyXMini0885.problem_phyx_mini_0885}
\uses{def:physics:phyx-mini-0885:phyxminiproblems-problemphyxmini0885-inductanceperlengthinmicrohenriespermeter, def:physics:phyx-mini-0885:phyxminiproblems-problemphyxmini0885-coaxialcablesetup, def:physics:phyx-mini-0885:phyxminiproblems-problemphyxmini0885-matcheswrittencoaxialcablescenario, def:physics:phyx-mini-0885:phyxminiproblems-problemphyxmini0885-matchessuppliedcoaxialcablefigure, def:physics:phyx-mini-0885:phyxminiproblems-problemphyxmini0885-matchesproblemradiusreadouts, def:physics:phyx-mini-0885:phyxminiproblems-problemphyxmini0885-hasphysicalcoaxialparameters, def:physics:phyx-mini-0885:phyxminiproblems-problemphyxmini0885-usestextbookvacuummagneticpermeability, def:physics:phyx-mini-0885:phyxminiproblems-problemphyxmini0885-satisfieshighfrequencycoaxialinductancelaw, def:physics:phyx-mini-0885:phyxminiproblems-problemphyxmini0885-answerchoice, def:physics:phyx-mini-0885:phyxminiproblems-problemphyxmini0885-recordeddatasetanswer, def:physics:phyx-mini-0885:phyxminiproblems-problemphyxmini0885-isroundedtodisplayedhundredth, def:physics:phyx-mini-0885:phyxminiproblems-problemphyxmini0885-answermatchesroundedinductance}
The assigned autoformalize agent should translate this physics problem into Lean declarations in the covered file, with theorem and lemma proof bodies written as `by sorry`.
\end{theorem}
\begin{proof}
This is an autoformalization task, not a proof task. Produce statements that can later be proved by the physics prover without weakening the physical model.
\end{proof}

% --- Archon declaration topology begin ---
\section{Declaration topology}
\begin{definition}[inductance Dimension]
  \label{def:physics:phyx-mini-0885:phyxminiproblems-problemphyxmini0885-inductancedimension}
  \lean{PhyXMiniProblems.ProblemPhyXMini0885.inductanceDimension}
  The physical dimension M L² C⁻² of electrical inductance
\end{definition}
\begin{proof}
  This is the declared model definition of inductance Dimension.
\end{proof}
\begin{definition}[inductance Per Length Dimension]
  \label{def:physics:phyx-mini-0885:phyxminiproblems-problemphyxmini0885-inductanceperlengthdimension}
  \lean{PhyXMiniProblems.ProblemPhyXMini0885.inductancePerLengthDimension}
  \uses{def:physics:phyx-mini-0885:phyxminiproblems-problemphyxmini0885-inductancedimension}
  The physical dimension M L C⁻² of inductance per length
\end{definition}
\begin{proof}
  This is the declared model definition of inductance Per Length Dimension.
\end{proof}
\begin{definition}[Length Quantity]
  \label{def:physics:phyx-mini-0885:phyxminiproblems-problemphyxmini0885-lengthquantity}
  \lean{PhyXMiniProblems.ProblemPhyXMini0885.LengthQuantity}
  A nonnegative, unit-independent physical length
\end{definition}
\begin{proof}
  This is the declared model definition of Length Quantity.
\end{proof}
\begin{definition}[Inductance Per Length Quantity]
  \label{def:physics:phyx-mini-0885:phyxminiproblems-problemphyxmini0885-inductanceperlengthquantity}
  \lean{PhyXMiniProblems.ProblemPhyXMini0885.InductancePerLengthQuantity}
  \uses{def:physics:phyx-mini-0885:phyxminiproblems-problemphyxmini0885-inductanceperlengthdimension}
  A nonnegative, unit-independent inductance per unit length
\end{definition}
\begin{proof}
  This is the declared model definition of Inductance Per Length Quantity.
\end{proof}
\begin{definition}[Magnetic Permeability Quantity]
  \label{def:physics:phyx-mini-0885:phyxminiproblems-problemphyxmini0885-magneticpermeabilityquantity}
  \lean{PhyXMiniProblems.ProblemPhyXMini0885.MagneticPermeabilityQuantity}
  \uses{def:physics:phyx-mini-0885:phyxminiproblems-problemphyxmini0885-inductanceperlengthdimension}
  A nonnegative, unit-independent magnetic permeability
\end{definition}
\begin{proof}
  This is the declared model definition of Magnetic Permeability Quantity.
\end{proof}
\begin{definition}[length Readout]
  \label{def:physics:phyx-mini-0885:phyxminiproblems-problemphyxmini0885-lengthreadout}
  \lean{PhyXMiniProblems.ProblemPhyXMini0885.lengthReadout}
  \uses{def:physics:phyx-mini-0885:phyxminiproblems-problemphyxmini0885-lengthquantity}
  Read a physical length using the length unit selected by units
\end{definition}
\begin{proof}
  This is the declared model definition of length Readout.
\end{proof}
\begin{definition}[length In Millimeters]
  \label{def:physics:phyx-mini-0885:phyxminiproblems-problemphyxmini0885-lengthinmillimeters}
  \lean{PhyXMiniProblems.ProblemPhyXMini0885.lengthInMillimeters}
  \uses{def:physics:phyx-mini-0885:phyxminiproblems-problemphyxmini0885-lengthquantity, def:physics:phyx-mini-0885:phyxminiproblems-problemphyxmini0885-lengthreadout}
  Read a physical length in millimetres
\end{definition}
\begin{proof}
  This is the declared model definition of length In Millimeters.
\end{proof}
\begin{definition}[inductance Per Length Readout]
  \label{def:physics:phyx-mini-0885:phyxminiproblems-problemphyxmini0885-inductanceperlengthreadout}
  \lean{PhyXMiniProblems.ProblemPhyXMini0885.inductancePerLengthReadout}
  \uses{def:physics:phyx-mini-0885:phyxminiproblems-problemphyxmini0885-inductanceperlengthquantity}
  Read an inductance per length in the coherent unit induced by units
\end{definition}
\begin{proof}
  This is the declared model definition of inductance Per Length Readout.
\end{proof}
\begin{definition}[inductance Per Length In Henries Per Meter]
  \label{def:physics:phyx-mini-0885:phyxminiproblems-problemphyxmini0885-inductanceperlengthinhenriespermeter}
  \lean{PhyXMiniProblems.ProblemPhyXMini0885.inductancePerLengthInHenriesPerMeter}
  \uses{def:physics:phyx-mini-0885:phyxminiproblems-problemphyxmini0885-inductanceperlengthquantity, def:physics:phyx-mini-0885:phyxminiproblems-problemphyxmini0885-inductanceperlengthreadout}
  Read an inductance per length in coherent-SI henries per metre
\end{definition}
\begin{proof}
  This is the declared model definition of inductance Per Length In Henries Per Meter.
\end{proof}
\begin{definition}[inductance Per Length In Microhenries Per Meter]
  \label{def:physics:phyx-mini-0885:phyxminiproblems-problemphyxmini0885-inductanceperlengthinmicrohenriespermeter}
  \lean{PhyXMiniProblems.ProblemPhyXMini0885.inductancePerLengthInMicrohenriesPerMeter}
  \uses{def:physics:phyx-mini-0885:phyxminiproblems-problemphyxmini0885-inductanceperlengthquantity, def:physics:phyx-mini-0885:phyxminiproblems-problemphyxmini0885-inductanceperlengthinhenriespermeter}
  Read an inductance per length in microhenries per metre
\end{definition}
\begin{proof}
  This is the declared model definition of inductance Per Length In Microhenries Per Meter.
\end{proof}
\begin{definition}[magnetic Permeability Readout]
  \label{def:physics:phyx-mini-0885:phyxminiproblems-problemphyxmini0885-magneticpermeabilityreadout}
  \lean{PhyXMiniProblems.ProblemPhyXMini0885.magneticPermeabilityReadout}
  \uses{def:physics:phyx-mini-0885:phyxminiproblems-problemphyxmini0885-magneticpermeabilityquantity}
  Read magnetic permeability in the coherent unit induced by units
\end{definition}
\begin{proof}
  This is the declared model definition of magnetic Permeability Readout.
\end{proof}
\begin{definition}[magnetic Permeability In Henries Per Meter]
  \label{def:physics:phyx-mini-0885:phyxminiproblems-problemphyxmini0885-magneticpermeabilityinhenriespermeter}
  \lean{PhyXMiniProblems.ProblemPhyXMini0885.magneticPermeabilityInHenriesPerMeter}
  \uses{def:physics:phyx-mini-0885:phyxminiproblems-problemphyxmini0885-magneticpermeabilityquantity, def:physics:phyx-mini-0885:phyxminiproblems-problemphyxmini0885-magneticpermeabilityreadout}
  Read magnetic permeability in coherent-SI henries per metre
\end{definition}
\begin{proof}
  This is the declared model definition of magnetic Permeability In Henries Per Meter.
\end{proof}
\begin{definition}[Cable Conductor]
  \label{def:physics:phyx-mini-0885:phyxminiproblems-problemphyxmini0885-cableconductor}
  \lean{PhyXMiniProblems.ProblemPhyXMini0885.CableConductor}
  The two cylindrical conductors identified in the scenario and image
\end{definition}
\begin{proof}
  This is the declared model definition of Cable Conductor.
\end{proof}
\begin{definition}[Conductor Model]
  \label{def:physics:phyx-mini-0885:phyxminiproblems-problemphyxmini0885-conductormodel}
  \lean{PhyXMiniProblems.ProblemPhyXMini0885.ConductorModel}
  The electromagnetic idealization assigned to a conductor
\end{definition}
\begin{proof}
  This is the declared model definition of Conductor Model.
\end{proof}
\begin{definition}[Insulating Material Model]
  \label{def:physics:phyx-mini-0885:phyxminiproblems-problemphyxmini0885-insulatingmaterialmodel}
  \lean{PhyXMiniProblems.ProblemPhyXMini0885.InsulatingMaterialModel}
  The two insulating material roles described in the written scenario
\end{definition}
\begin{proof}
  This is the declared model definition of Insulating Material Model.
\end{proof}
\begin{definition}[Magnetic Material Model]
  \label{def:physics:phyx-mini-0885:phyxminiproblems-problemphyxmini0885-magneticmaterialmodel}
  \lean{PhyXMiniProblems.ProblemPhyXMini0885.MagneticMaterialModel}
  Magnetic behavior of the dielectric fill used by the textbook formula
\end{definition}
\begin{proof}
  This is the declared model definition of Magnetic Material Model.
\end{proof}
\begin{definition}[Signal Frequency Regime]
  \label{def:physics:phyx-mini-0885:phyxminiproblems-problemphyxmini0885-signalfrequencyregime}
  \lean{PhyXMiniProblems.ProblemPhyXMini0885.SignalFrequencyRegime}
  Frequency regimes distinguished by the cable model
\end{definition}
\begin{proof}
  This is the declared model definition of Signal Frequency Regime.
\end{proof}
\begin{definition}[Figure Label]
  \label{def:physics:phyx-mini-0885:phyxminiproblems-problemphyxmini0885-figurelabel}
  \lean{PhyXMiniProblems.ProblemPhyXMini0885.FigureLabel}
  Text or symbolic annotations visible in the supplied cable image
\end{definition}
\begin{proof}
  This is the declared model definition of Figure Label.
\end{proof}
\begin{definition}[Coaxial Cable Figure]
  \label{def:physics:phyx-mini-0885:phyxminiproblems-problemphyxmini0885-coaxialcablefigure}
  \lean{PhyXMiniProblems.ProblemPhyXMini0885.CoaxialCableFigure}
  \uses{def:physics:phyx-mini-0885:phyxminiproblems-problemphyxmini0885-cableconductor, def:physics:phyx-mini-0885:phyxminiproblems-problemphyxmini0885-figurelabel}
  Literal presentation features of image 885.png. These fields describe the two concentric cylinders, dashed common axis, and the arrows labelled r₁ and r₂; they contain no inductance value
\end{definition}
\begin{proof}
  This is the declared model definition of Coaxial Cable Figure.
\end{proof}
\begin{definition}[Coaxial Cable Setup]
  \label{def:physics:phyx-mini-0885:phyxminiproblems-problemphyxmini0885-coaxialcablesetup}
  \lean{PhyXMiniProblems.ProblemPhyXMini0885.CoaxialCableSetup}
  \uses{def:physics:phyx-mini-0885:phyxminiproblems-problemphyxmini0885-lengthquantity, def:physics:phyx-mini-0885:phyxminiproblems-problemphyxmini0885-inductanceperlengthquantity, def:physics:phyx-mini-0885:phyxminiproblems-problemphyxmini0885-magneticpermeabilityquantity, def:physics:phyx-mini-0885:phyxminiproblems-problemphyxmini0885-cableconductor, def:physics:phyx-mini-0885:phyxminiproblems-problemphyxmini0885-conductormodel, def:physics:phyx-mini-0885:phyxminiproblems-problemphyxmini0885-insulatingmaterialmodel, def:physics:phyx-mini-0885:phyxminiproblems-problemphyxmini0885-magneticmaterialmodel, def:physics:phyx-mini-0885:phyxminiproblems-problemphyxmini0885-signalfrequencyregime, def:physics:phyx-mini-0885:phyxminiproblems-problemphyxmini0885-coaxialcablefigure}
  Independent physical data for the cable. The requested inductance per length is an unconstrained physical quantity here; only the governing law below relates it to the radii and permeability
\end{definition}
\begin{proof}
  This is the declared model definition of Coaxial Cable Setup.
\end{proof}
\begin{definition}[Matches Written Coaxial Cable Scenario]
  \label{def:physics:phyx-mini-0885:phyxminiproblems-problemphyxmini0885-matcheswrittencoaxialcablescenario}
  \lean{PhyXMiniProblems.ProblemPhyXMini0885.MatchesWrittenCoaxialCableScenario}
  \uses{def:physics:phyx-mini-0885:phyxminiproblems-problemphyxmini0885-coaxialcablesetup}
  The qualitative cable model stated in the written problem
\end{definition}
\begin{proof}
  This is the declared model definition of Matches Written Coaxial Cable Scenario.
\end{proof}
\begin{definition}[Matches Supplied Coaxial Cable Figure]
  \label{def:physics:phyx-mini-0885:phyxminiproblems-problemphyxmini0885-matchessuppliedcoaxialcablefigure}
  \lean{PhyXMiniProblems.ProblemPhyXMini0885.MatchesSuppliedCoaxialCableFigure}
  \uses{def:physics:phyx-mini-0885:phyxminiproblems-problemphyxmini0885-coaxialcablesetup}
  Primary-raster evidence. It records the labelled coaxial geometry only; the radius calibrations come from the question text and are stated separately
\end{definition}
\begin{proof}
  This is the declared model definition of Matches Supplied Coaxial Cable Figure.
\end{proof}
\begin{definition}[Matches Problem Radius Readouts]
  \label{def:physics:phyx-mini-0885:phyxminiproblems-problemphyxmini0885-matchesproblemradiusreadouts}
  \lean{PhyXMiniProblems.ProblemPhyXMini0885.MatchesProblemRadiusReadouts}
  \uses{def:physics:phyx-mini-0885:phyxminiproblems-problemphyxmini0885-lengthinmillimeters, def:physics:phyx-mini-0885:phyxminiproblems-problemphyxmini0885-coaxialcablesetup}
  The two numerical radius readouts supplied in the question. They constrain only the cable geometry and contain no inductance value
\end{definition}
\begin{proof}
  This is the declared model definition of Matches Problem Radius Readouts.
\end{proof}
\begin{definition}[Has Physical Coaxial Parameters]
  \label{def:physics:phyx-mini-0885:phyxminiproblems-problemphyxmini0885-hasphysicalcoaxialparameters}
  \lean{PhyXMiniProblems.ProblemPhyXMini0885.HasPhysicalCoaxialParameters}
  \uses{def:physics:phyx-mini-0885:phyxminiproblems-problemphyxmini0885-lengthreadout, def:physics:phyx-mini-0885:phyxminiproblems-problemphyxmini0885-magneticpermeabilityinhenriespermeter, def:physics:phyx-mini-0885:phyxminiproblems-problemphyxmini0885-coaxialcablesetup}
  Positivity and nesting required by the logarithmic coaxial formula
\end{definition}
\begin{proof}
  This is the declared model definition of Has Physical Coaxial Parameters.
\end{proof}
\begin{definition}[Uses Textbook Vacuum Magnetic Permeability]
  \label{def:physics:phyx-mini-0885:phyxminiproblems-problemphyxmini0885-usestextbookvacuummagneticpermeability}
  \lean{PhyXMiniProblems.ProblemPhyXMini0885.UsesTextbookVacuumMagneticPermeability}
  \uses{def:physics:phyx-mini-0885:phyxminiproblems-problemphyxmini0885-magneticpermeabilityinhenriespermeter, def:physics:phyx-mini-0885:phyxminiproblems-problemphyxmini0885-coaxialcablesetup}
  The nonmagnetic fill has the permeability of free space. Physlib's Electromagnetism.EMSystem.μ₀ supplies the electromagnetic-system parameter; the second field gives the coherent-SI calibration μ₀ = 4π × 10⁻⁷ H/m. Neither field mentions the requested inductance
\end{definition}
\begin{proof}
  This is the declared model definition of Uses Textbook Vacuum Magnetic Permeability.
\end{proof}
\begin{definition}[Satisfies High Frequency Coaxial Inductance Law]
  \label{def:physics:phyx-mini-0885:phyxminiproblems-problemphyxmini0885-satisfieshighfrequencycoaxialinductancelaw}
  \lean{PhyXMiniProblems.ProblemPhyXMini0885.SatisfiesHighFrequencyCoaxialInductanceLaw}
  \uses{def:physics:phyx-mini-0885:phyxminiproblems-problemphyxmini0885-lengthreadout, def:physics:phyx-mini-0885:phyxminiproblems-problemphyxmini0885-inductanceperlengthreadout, def:physics:phyx-mini-0885:phyxminiproblems-problemphyxmini0885-magneticpermeabilityreadout, def:physics:phyx-mini-0885:phyxminiproblems-problemphyxmini0885-coaxialcablesetup}
  The general governing high-frequency coaxial-cable law L' = μ / (2π) * log (r₂ / r₁). It is stated in every coherent unit system. The radius ratio is dimensionless, and magnetic permeability and inductance per length have the same dimension. This law contains no answer choice or rounded target value
\end{definition}
\begin{proof}
  This is the declared model definition of Satisfies High Frequency Coaxial Inductance Law.
\end{proof}
\begin{definition}[Answer Choice]
  \label{def:physics:phyx-mini-0885:phyxminiproblems-problemphyxmini0885-answerchoice}
  \lean{PhyXMiniProblems.ProblemPhyXMini0885.AnswerChoice}
  The four answer labels shown with the problem
\end{definition}
\begin{proof}
  This is the declared model definition of Answer Choice.
\end{proof}
\begin{definition}[displayed Inductance In Microhenries Per Meter]
  \label{def:physics:phyx-mini-0885:phyxminiproblems-problemphyxmini0885-answerchoice-displayedinductanceinmicrohenriespermeter}
  \lean{PhyXMiniProblems.ProblemPhyXMini0885.AnswerChoice.displayedInductanceInMicrohenriesPerMeter}
  \uses{def:physics:phyx-mini-0885:phyxminiproblems-problemphyxmini0885-answerchoice}
  The displayed numerical readout for each choice, in μH/m
\end{definition}
\begin{proof}
  This is the declared model definition of displayed Inductance In Microhenries Per Meter.
\end{proof}
\begin{definition}[recorded Dataset Answer]
  \label{def:physics:phyx-mini-0885:phyxminiproblems-problemphyxmini0885-recordeddatasetanswer}
  \lean{PhyXMiniProblems.ProblemPhyXMini0885.recordedDatasetAnswer}
  \uses{def:physics:phyx-mini-0885:phyxminiproblems-problemphyxmini0885-answerchoice}
  The answer label recorded by the source dataset
\end{definition}
\begin{proof}
  This is the declared model definition of recorded Dataset Answer.
\end{proof}
\begin{definition}[Is Rounded To Displayed Hundredth]
  \label{def:physics:phyx-mini-0885:phyxminiproblems-problemphyxmini0885-isroundedtodisplayedhundredth}
  \lean{PhyXMiniProblems.ProblemPhyXMini0885.IsRoundedToDisplayedHundredth}
  displayed is the nearest-hundredth presentation of actual when their difference is strictly less than half of one hundredth
\end{definition}
\begin{proof}
  This is the declared model definition of Is Rounded To Displayed Hundredth.
\end{proof}
\begin{definition}[Answer Matches Rounded Inductance]
  \label{def:physics:phyx-mini-0885:phyxminiproblems-problemphyxmini0885-answermatchesroundedinductance}
  \lean{PhyXMiniProblems.ProblemPhyXMini0885.AnswerMatchesRoundedInductance}
  \uses{def:physics:phyx-mini-0885:phyxminiproblems-problemphyxmini0885-inductanceperlengthinmicrohenriespermeter, def:physics:phyx-mini-0885:phyxminiproblems-problemphyxmini0885-coaxialcablesetup, def:physics:phyx-mini-0885:phyxminiproblems-problemphyxmini0885-answerchoice, def:physics:phyx-mini-0885:phyxminiproblems-problemphyxmini0885-isroundedtodisplayedhundredth}
  A choice agrees with the cable's computed inductance after rounding
\end{definition}
\begin{proof}
  This is the declared model definition of Answer Matches Rounded Inductance.
\end{proof}
% --- Archon declaration topology end ---
% --- Archon physics formalization source end ---
